% ═══════════════════════════════════════════════════════════════
% LB-01a: Lehrbuchseite - Grundgesetz der Mechanik (F = m · a)
% Physik Klasse 10 MR
% ═══════════════════════════════════════════════════════════════

\documentclass[11pt, a4paper]{article}

\usepackage[utf8]{inputenc}
\usepackage[T1]{fontenc}
\usepackage[ngerman]{babel}

\usepackage{helvet}
\renewcommand{\familydefault}{\sfdefault}

\usepackage{microtype}
\usepackage[margin=1.8cm, top=2cm, bottom=1.5cm]{geometry}
\usepackage{enumitem}
\usepackage{tabularx}
\usepackage{booktabs}
\usepackage{array}
\usepackage{multicol}
\usepackage{tikz}
\usetikzlibrary{calc, arrows.meta, decorations.pathreplacing, positioning}

\usepackage{amsmath, amssymb}
\usepackage{siunitx}
\sisetup{locale = DE, per-mode = symbol}

\usepackage{fancyhdr}
\usepackage{lastpage}

\usepackage{xcolor}
\usepackage{tcolorbox}
\tcbuselibrary{skins, breakable}

% ═══════════════════════════════════════════════════════════════
% FARBEN
% ═══════════════════════════════════════════════════════════════

\definecolor{physik}{HTML}{2c5aa0}
\definecolor{hellgrau}{HTML}{F5F5F5}
\definecolor{mittelgrau}{HTML}{888888}
\definecolor{rahmen}{HTML}{CCCCCC}
\definecolor{akzent}{HTML}{e8f0f8}
\definecolor{erkenntnisfarbe}{HTML}{fff3e0}

\colorlet{fachfarbe}{physik}

% ═══════════════════════════════════════════════════════════════
% VARIABLEN
% ═══════════════════════════════════════════════════════════════

\newcommand{\fach}{Physik}
\newcommand{\thema}{Grundgesetz der Mechanik}
\newcommand{\klasse}{10 MR}

% ═══════════════════════════════════════════════════════════════
% KOPF- UND FUSSZEILE
% ═══════════════════════════════════════════════════════════════

\pagestyle{fancy}
\fancyhf{}
\renewcommand{\headrulewidth}{0.4pt}
\renewcommand{\footrulewidth}{0pt}
\lhead{\small \fach{} Klasse \klasse{} | \thema{}}
\rhead{\small Lehrbuchseite}
\cfoot{\small\textcolor{mittelgrau}{Seite \thepage{} von \pageref{LastPage}}}

% ═══════════════════════════════════════════════════════════════
% BOXEN
% ═══════════════════════════════════════════════════════════════

\newtcolorbox{definitionsbox}[1][]{
    colback=hellgrau,
    colframe=hellgrau,
    boxrule=0pt,
    arc=0pt,
    left=10pt, right=10pt, top=10pt, bottom=8pt,
    borderline west={3pt}{0pt}{fachfarbe},
    before skip=8pt, after skip=8pt,
    #1
}

\newtcolorbox{erkenntnisbox}[1][]{
    colback=erkenntnisfarbe,
    colframe=erkenntnisfarbe,
    boxrule=0pt,
    arc=0pt,
    left=10pt, right=10pt, top=8pt, bottom=8pt,
    borderline west={3pt}{0pt}{orange!80!black},
    before skip=6pt, after skip=6pt,
    #1
}

\newtcolorbox{formelbox}{
    colback=white,
    colframe=fachfarbe,
    boxrule=2pt,
    arc=0pt,
    left=10pt, right=10pt, top=10pt, bottom=10pt,
    before skip=8pt, after skip=8pt
}

\newtcolorbox{beispielbox}[1][]{
    colback=akzent,
    colframe=akzent,
    boxrule=0pt,
    arc=0pt,
    left=10pt, right=10pt, top=8pt, bottom=6pt,
    borderline north={2pt}{0pt}{fachfarbe},
    before skip=6pt, after skip=6pt,
    #1
}

\newtcolorbox{merksatzbox}{
    colback=white,
    colframe=fachfarbe,
    boxrule=1pt,
    arc=0pt,
    left=10pt, right=10pt, top=8pt, bottom=8pt,
    before skip=6pt, after skip=6pt
}

% ═══════════════════════════════════════════════════════════════
% BEFEHLE
% ═══════════════════════════════════════════════════════════════

\newcommand{\begriff}[1]{\textbf{\textcolor{fachfarbe}{#1}}}
\newcommand{\einheit}[1]{\textcolor{mittelgrau}{[#1]}}

\setlength{\parindent}{0pt}
\setlength{\parskip}{6pt}

% ═══════════════════════════════════════════════════════════════
% DOKUMENT
% ═══════════════════════════════════════════════════════════════

\begin{document}

% ═══════════════════════════════════════════════════════════════
% HEADER
% ═══════════════════════════════════════════════════════════════

\begin{center}
    {\LARGE\bfseries\textcolor{fachfarbe}{Das Grundgesetz der Mechanik}}\\[0.3em]
    {\large $F = m \cdot a$ \quad (Newtons 2. Gesetz)}
\end{center}
\vspace{-0.2em}
\hrule
\vspace{0.8em}

\begin{multicols}{2}

% ═══════════════════════════════════════════════════════════════
% MOTIVATION: WARUM IST DAS WICHTIG?
% ═══════════════════════════════════════════════════════════════

\begin{tcolorbox}[
    colback=fachfarbe!10,
    colframe=fachfarbe,
    boxrule=1.5pt,
    arc=0pt,
    left=8pt, right=8pt, top=6pt, bottom=6pt,
    title={\textcolor{white}{\textbf{Warum ist das wichtig für DICH?}}}
]
\small
\textbf{Denk mal nach:}
\begin{itemize}[leftmargin=1.2em, itemsep=2pt]
    \item Warum beschleunigt ein \textbf{Sportwagen} schneller als ein LKW?
    \item Warum ist ein voller \textbf{Einkaufswagen} so schwer zu schieben?
    \item Wie berechnen Ingenieure die \textbf{Bremswege} für deine Sicherheit?
\end{itemize}

\vspace{0.3em}
\textbf{AHA-Fakt:} Du erlebst das jeden Tag! Je stärker du beim Fahrradfahren in die Pedale trittst (mehr Kraft), desto schneller wirst du (mehr Beschleunigung). Mit vollen Taschen am Lenker (mehr Masse) wird das Beschleunigen schwerer!
\end{tcolorbox}

\vspace{0.4em}

% ═══════════════════════════════════════════════════════════════
% 1. RESULTIERENDE KRAFT
% ═══════════════════════════════════════════════════════════════

\section*{\textcolor{fachfarbe}{1. Resultierende Kraft}}
{\small\textcolor{mittelgrau}{$\rightarrow$ AB Kasten 1}}

Oft wirken \textbf{mehrere Kräfte} gleichzeitig auf einen Körper. Die \begriff{resultierende Kraft} ist die Summe aller Kräfte.

\begin{definitionsbox}
\textbf{Gleiche Richtung:} Kräfte addieren\\
$F_{\text{res}} = F_1 + F_2$

\vspace{0.3em}
\textbf{Entgegengesetzt:} Kräfte subtrahieren\\
$F_{\text{res}} = F_1 - F_2$
\end{definitionsbox}

\begin{center}
\begin{tikzpicture}[scale=0.6]
    % Seil
    \draw[line width=4pt, brown!60!black] (-3.3,0.5) -- (3.3,0.5);

    % Team links - 3 Strichmännchen
    \foreach \offset in {0, 0.9, 1.8} {
        % Körper
        \draw[line width=1.5pt] (-3.5-\offset,0.5) -- (-3.8-\offset,1.0); % Arm zum Seil
        \draw[line width=1.5pt] (-3.8-\offset,1.0) -- (-3.8-\offset,1.8); % Körper
        \draw[fill=white, draw=black, line width=1pt] (-3.8-\offset,2.0) circle (0.22); % Kopf
        \draw[line width=1.5pt] (-3.8-\offset,1.0) -- (-4.1-\offset,0); % Bein hinten
        \draw[line width=1.5pt] (-3.8-\offset,1.0) -- (-3.5-\offset,0); % Bein vorne
    }

    % Team rechts - 3 Strichmännchen
    \foreach \offset in {0, 0.9, 1.8} {
        % Körper
        \draw[line width=1.5pt] (3.5+\offset,0.5) -- (3.8+\offset,1.0); % Arm zum Seil
        \draw[line width=1.5pt] (3.8+\offset,1.0) -- (3.8+\offset,1.8); % Körper
        \draw[fill=white, draw=black, line width=1pt] (3.8+\offset,2.0) circle (0.22); % Kopf
        \draw[line width=1.5pt] (3.8+\offset,1.0) -- (4.1+\offset,0); % Bein hinten
        \draw[line width=1.5pt] (3.8+\offset,1.0) -- (3.5+\offset,0); % Bein vorne
    }

    % Angriffspunkte am Seil (x)
    \node[font=\bfseries] at (-3.3,0.5) {\textcolor{blue!70!black}{$\times$}};
    \node[font=\bfseries] at (3.3,0.5) {\textcolor{red!70!black}{$\times$}};

    % Kraftpfeil links (am Seil-Ende angreifend)
    \draw[-{Stealth[length=4mm]}, blue!70!black, line width=2.5pt] (-3.3,0.5) -- (-5.5,0.5);
    \node[blue!70!black, font=\bfseries] at (-4.4,1.0) {\small $F_{\text{links}}$};
    \node[blue!70!black] at (-4.4,-0.2) {\small 800 N};

    % Kraftpfeil rechts (am Seil-Ende angreifend)
    \draw[-{Stealth[length=4mm]}, red!70!black, line width=2.5pt] (3.3,0.5) -- (5.5,0.5);
    \node[red!70!black, font=\bfseries] at (4.4,1.0) {\small $F_{\text{rechts}}$};
    \node[red!70!black] at (4.4,-0.2) {\small 900 N};

    % Resultierende Kraft (Mitte, nach rechts)
    \draw[-{Stealth[length=4mm]}, fachfarbe, line width=3pt] (0,2.0) -- (1.8,2.0);
    \node[fachfarbe, font=\bfseries] at (0.9,2.5) {\small $F_{\text{res}} = 100\,$N $\rightarrow$};
\end{tikzpicture}
\end{center}

% ═══════════════════════════════════════════════════════════════
% 2a. ERKENNTNIS KRAFT
% ═══════════════════════════════════════════════════════════════

\section*{\textcolor{fachfarbe}{2a. Wie wirkt Kraft?}}
{\small\textcolor{mittelgrau}{$\rightarrow$ AB Kasten 2a}}

Zwei Kisten mit \textbf{gleicher Masse} werden mit \textbf{unterschiedlicher Kraft} geschoben.

\begin{center}
\begin{tikzpicture}[scale=0.6]
    % Obere Kiste (kleine Kraft) - 50% breiter: 1.2 → 1.8
    \draw[fill=blue!20, draw=fachfarbe, line width=1.5pt] (0,1.4) rectangle (1.8,2.4);
    \node at (0.9,1.9) {\small\textbf{10 kg}};
    \draw[-{Stealth[length=3mm]}, blue!70!black, line width=2pt] (2.0,1.9) -- (3.0,1.9);
    \node[blue!70!black] at (3.7,1.9) {\small 20 N};

    % Untere Kiste (große Kraft) - 50% breiter: 1.2 → 1.8
    \draw[fill=red!20, draw=fachfarbe, line width=1.5pt] (0,0) rectangle (1.8,1.0);
    \node at (0.9,0.5) {\small\textbf{10 kg}};
    \draw[-{Stealth[length=3mm]}, red!70!black, line width=2.5pt] (2.0,0.5) -- (4.0,0.5);
    \node[red!70!black] at (4.7,0.5) {\small 50 N};
\end{tikzpicture}
\end{center}

\begin{erkenntnisbox}
\textbf{Erkenntnis 1:}\\
Je größer die \textbf{Kraft}, desto \textbf{größer} die Beschleunigung.
\end{erkenntnisbox}

% ═══════════════════════════════════════════════════════════════
% 2b. ERKENNTNIS MASSE
% ═══════════════════════════════════════════════════════════════

\section*{\textcolor{fachfarbe}{2b. Wie wirkt Masse?}}
{\small\textcolor{mittelgrau}{$\rightarrow$ AB Kasten 2b}}

Zwei Einkaufswagen mit \textbf{unterschiedlicher Masse} werden mit \textbf{gleicher Kraft} geschoben.

\begin{center}
\begin{tikzpicture}[scale=0.6]
    % Oberer Wagen (leicht) - 50% breiter: 1.0 → 1.5
    \draw[fill=green!20, draw=fachfarbe, line width=1.5pt] (0,1.4) rectangle (1.5,2.3);
    \node at (0.75,1.85) {\small\textbf{10 kg}};
    \draw[-{Stealth[length=3mm]}, fachfarbe, line width=2pt] (1.7,1.85) -- (3.2,1.85);
    \node[fachfarbe] at (3.9,1.85) {\small 50 N};

    % Unterer Wagen (schwer) - 50% breiter: 1.5 → 2.25
    \draw[fill=red!20, draw=fachfarbe, line width=1.5pt] (0,0) rectangle (2.25,1.1);
    \node at (1.125,0.55) {\small\textbf{50 kg}};
    \draw[-{Stealth[length=3mm]}, fachfarbe, line width=2pt] (2.45,0.55) -- (3.95,0.55);
    \node[fachfarbe] at (4.65,0.55) {\small 50 N};
\end{tikzpicture}
\end{center}

\begin{erkenntnisbox}
\textbf{Erkenntnis 2:}\\
Je größer die \textbf{Masse}, desto \textbf{kleiner} die Beschleunigung.
\end{erkenntnisbox}

% ═══════════════════════════════════════════════════════════════
% 2c. VON DER BEOBACHTUNG ZUR FORMEL
% ═══════════════════════════════════════════════════════════════

\vspace{0.3em}
\begin{definitionsbox}
\textbf{Von der Beobachtung zur Formel:}

\vspace{0.2em}
\textbf{Erkenntn. 1:} $a \sim F$ {\small(proportional)}

\textbf{Erkenntn. 2:} $a \sim \frac{1}{m}$ {\small(umgek. prop.)}

\vspace{0.2em}
\hrule
\vspace{0.2em}
Zusammen: $a = \frac{F}{m}$ \enspace$\Rightarrow$\enspace $\mathbf{F = m \cdot a}$
\end{definitionsbox}
{\small\textcolor{mittelgrau}{$\rightarrow$ AB Kasten 2c}}

\columnbreak

% ═══════════════════════════════════════════════════════════════
% 3. GRUNDGESETZ DER MECHANIK
% ═══════════════════════════════════════════════════════════════

\section*{\textcolor{fachfarbe}{3. Das Grundgesetz}}
{\small\textcolor{mittelgrau}{$\rightarrow$ AB Kasten 3}}

Newton fasste beide Erkenntnisse in einer Formel zusammen:

\begin{formelbox}
\begin{center}
{\Huge $F = m \cdot a$}

\vspace{0.6em}
\begin{tabular}{lll}
$F$ & = Kraft & \einheit{N} \\[0.3em]
$m$ & = Masse & \einheit{kg} \\[0.3em]
$a$ & = Beschleunigung & \einheit{m/s²}
\end{tabular}
\end{center}
\end{formelbox}

\begin{merksatzbox}
\textbf{Merksatz:}\\
Je größer die Kraft $\rightarrow$ größere Beschleunigung\\
Je größer die Masse $\rightarrow$ kleinere Beschleunigung
\end{merksatzbox}

% ═══════════════════════════════════════════════════════════════
% 4. FORMEL-DREIECK
% ═══════════════════════════════════════════════════════════════

\section*{\textcolor{fachfarbe}{4. Das Formel-Dreieck}}
{\small\textcolor{mittelgrau}{$\rightarrow$ AB Kasten 4}}

Das Formel-Dreieck ist ein \textbf{Trick}, um Formeln schnell umzustellen:

\vspace{0.2em}
\textbf{So geht's:} Halte den Finger auf die \textbf{gesuchte Größe}. Was übrig bleibt, ist die Formel!

\vspace{0.3em}

\begin{center}
\begin{tikzpicture}[scale=1.0]
    % Dreieck
    \draw[fachfarbe, line width=2pt] (0,0) -- (3.2,0) -- (1.6,2.4) -- cycle;
    % Horizontale Linie
    \draw[fachfarbe, line width=1.5pt] (0.45,1.2) -- (2.75,1.2);
    % Vertikale Linie
    \draw[fachfarbe, line width=1.5pt] (1.6,0) -- (1.6,1.2);
    % Beschriftung
    \node at (1.6,1.75) {\LARGE\textbf{F}};
    \node at (0.8,0.5) {\LARGE\textbf{m}};
    \node at (2.4,0.5) {\LARGE\textbf{a}};
\end{tikzpicture}
\end{center}

\vspace{0.3em}

\textbf{Anleitung:} Decke die gesuchte Größe ab!

\begin{tabular}{@{}lll@{}}
\textbf{Gesucht} & \textbf{Abdecken} & \textbf{Formel} \\[0.3em]
Kraft & F & $F = m \cdot a$ \\[0.2em]
Masse & m & $m = \dfrac{F}{a}$ \\[0.5em]
Beschleunigung & a & $a = \dfrac{F}{m}$
\end{tabular}

\vspace{0.5em}

\begin{beispielbox}
\textbf{Beispiel:} $m = 5\,$kg, $a = 2\,$m/s²

$F = m \cdot a = 5\,\text{kg} \cdot 2\,\text{m/s}^2 = \mathbf{10\,N}$
\end{beispielbox}

\end{multicols}

\vspace{0.3em}
\hrule
\vspace{0.5em}

% ═══════════════════════════════════════════════════════════════
% ZUSAMMENFASSUNG
% ═══════════════════════════════════════════════════════════════

\begin{merksatzbox}
\begin{center}
\textbf{Zusammenfassung: Das Wichtigste auf einen Blick}
\end{center}

\begin{multicols}{3}
\begin{center}
\textbf{Kraft berechnen}\\[0.3em]
$F = m \cdot a$\\[0.2em]
\small Masse $\times$ Beschleunigung
\end{center}

\columnbreak

\begin{center}
\textbf{Masse berechnen}\\[0.3em]
$m = \dfrac{F}{a}$\\[0.2em]
\small Kraft / Beschleunigung
\end{center}

\columnbreak

\begin{center}
\textbf{Beschleunigung berechnen}\\[0.3em]
$a = \dfrac{F}{m}$\\[0.2em]
\small Kraft / Masse
\end{center}
\end{multicols}
\end{merksatzbox}

% ═══════════════════════════════════════════════════════════════
% BEISPIELRECHNUNGEN
% ═══════════════════════════════════════════════════════════════

\vspace{0.3em}

\begin{beispielbox}
\textbf{Beispielrechnungen zum Üben:}

\begin{multicols}{3}
\textbf{1. Kraft gesucht}\\
$m = 20\,$kg, $a = 5\,$m/s²\\
$F = 20 \cdot 5 = \mathbf{100\,N}$

\columnbreak

\textbf{2. Masse gesucht}\\
$F = 200\,$N, $a = 10\,$m/s²\\
$m = 200 \div 10 = \mathbf{20\,kg}$

\columnbreak

\textbf{3. Beschleunigung gesucht}\\
$F = 1000\,$N, $m = 200\,$kg\\
$a = 1000 \div 200 = \mathbf{5\,m/s^2}$
\end{multicols}
\end{beispielbox}

\vspace{0.5em}
\begin{center}
\fcolorbox{fachfarbe}{fachfarbe!10}{%
\parbox{0.9\linewidth}{%
\centering\small
\textbf{Jetzt bist du dran!} $\rightarrow$ \textbf{AB Aufgaben 1--5} bearbeiten
}%
}
\end{center}

\vfill

\begin{center}
\small\textcolor{mittelgrau}{Physik Klasse 10 MR | Grundgesetz der Mechanik}
\end{center}

\end{document}
